\subsection{\texorpdfstring{Divisibility of polynomials in one indeterminate \(^{1}\)}{Divisibility of polynomials in one indeterminate [1]}}

PROPOSITION 11. --- Let K be a commutative field.

\begin{enumerate}
\def\labelenumi{(\roman{enumi})}
\item
  For every non-zero ideal \(\mathfrak{a}\) of \(\mathbf{K}[\mathbf{X}]\) there exists precisely one monic polynomial \(f\) in \(\mathbf{K}[\mathbf{X}]\) such that \(\mathfrak{a} = (f)\) .
\item
  Let \(f_{1}\) and \(f_{2}\) be in \(\mathbf{K}[\mathbf{X}]\) ; for \((f_{1}) = (f_{2})\) to hold it is necessary and sufficient that there exist a non-zero element \(\lambda\) of \(\mathbf{K}\) such that \(f_{2} = \lambda f_{1}\) .
\end{enumerate}

Let us prove (ii), the sufficiency of the stated condition being clear. The case where \(f_{1}\) and \(f_{2}\) generate the zero ideal is trivial. Thus assume that the non-zero polynomials \(f_{1}\) and \(f_{2}\) generate the same ideal of K{[}X{]}. Then there exist polynomials \(u_{1}\) and \(u_{2}\) such that \(f_{1}=u_{1}f_{2}\) and \(f_{2}=u_{2}f_{1}\) ; it follows that \(u_{1}u_{2} = 1\) , whence \(\deg u_{1} + \deg u_{2} = 0\) and so \(\deg u_{2} = 0\) . We have thus shown that \(u_{2}\) is a non-zero element of \(\mathbf{K}\) .

To prove (i), let \(f\) be a monic polynomial in \(\mathfrak{a}\) of least possible degree. Given \(g\) in \(\mathfrak{a}\) , let \(u\) and \(v\) be the quotient and remainder of the Euclidean division of \(g\) by \(f\) ; then \(v = g - uf\) belongs to \(\mathfrak{a}\) and we have \(\deg v < \deg f\) ; if \(v\) were non-zero, there would be a non-zero element \(\lambda\) of \(K\) such that \(\lambda v\) is monic, and since \(\lambda v \in \mathfrak{a}\) , this would contradict the definition of \(f\) . We thus have \(\mathfrak{a} = (f)\) ; the uniqueness of the monic polynomial \(f\) such that \(\mathfrak{a} = (f)\) now follows from (ii).

PROPOSITION 12. --- Let K be a commutative field and f, g two elements of K{[}X{]}. For every polynomial d in K{[}X{]} the following properties are equivalent:

\begin{enumerate}
\def\labelenumi{(\roman{enumi})}
\item
  The polynomial \(d\) divides \(f\) and \(g\) and every polynomial which divides both \(f\) and \(g\) divides \(d\) .
\item
  The polynomial \(d\) divides \(f\) and \(g\) and there exist two polynomials \(u\) and \(v\) such that \(d = uf + vg\) .
\item
  The relation \((d) = (f) + (g)\) holds between ideals in \(\mathbf{K}[\mathbf{X}]\) .
\end{enumerate}

The polynomial d is determined up to multiplication by a non-zero element of K by these properties. If f and g are not both zero, then \(d \neq 0\) and the degree of d majorizes the degree of every polynomial dividing both f and g.

When \(f\) and \(g\) are zero, each of the properties (i) to (iii) is satisfied only for \(d = 0\) , hence they are then equivalent. Henceforth we assume that \(f, g\) are not both 0 and we denote by \(\mathfrak{a}\) the ideal \((f) + (g)\) of \(\mathbf{K}[\mathbf{X}]\) .

We remark that for any polynomials \(u\) and \(v\) in \(\mathbf{K}[\mathbf{X}]\) the properties \((u) \supset (v)\) and \(\langle u \text{ divides } v \rangle\) are equivalent. The assertion (ii) is thus equivalent to \(\langle (d) \supset (f) \text{ and } (d) \supset (g) \text{ and } d \in (f) + (g) \rangle\) , that is (iii). It is clear that (ii) implies (i). Finally suppose that (i) holds; we have \((d) \supset (f)\) and \((d) \supset (g)\) , whence \((d) \supset \mathfrak{a}\) ; on the other hand, by Prop. 11 (IV, p.~12) there exists a polynomial \(d_1\) such that \(\mathfrak{a} = (d_1)\) ; since \(d_1\) divides both \(f\) and \(g\) , it divides \(d\) by hypothesis, whence \((d) \subset \mathfrak{a}\) , and finally we have \((d) = \mathfrak{a}\) , that is, (iii).

The other assertions of Prop. 12 are immediate consequences of Prop. 11 applied to the ideal \(\mathfrak{a} = (f) + (g)\) .

DEFINITION 1. --- With the notation of Prop. 12 we say that d is a greatest common divisor (gcd for short) of f and g. We say that f and g are relatively prime or that f is prime to g if 1 is a gcd of f and g.

To say that f and g are relatively prime thus means that there exist polynomials u and v in K{[}X{]} such that \(uf + vg = 1\) .

COROLLARY 1. --- Let \(d\) be a gcd of \(f\) and \(g\) and \(\mathbf{K}'\) a commutative field containing \(\mathbf{K}\) as subfield. Then \(d\) is a gcd of \(f\) and \(g\) considered as elements of \(\mathbf{K}'[\mathbf{X}]\) .

This follows from Prop. 12, (iii).

COROLLARY 2. --- Let d be a gcd of f and g.

\begin{enumerate}
\def\labelenumi{(\roman{enumi})}
\item
  If \(u \in \mathbf{K}[\mathbf{X}]\) , du is a gcd of fu and gu.
\item
  If \(v \in \mathbf{K}[\mathbf{X}]\) is a divisor ( \(\neq 0\) ) of \(f\) and \(g\) , then \(d / v\) is a gcd of \(f / v\) and \(g / v\) .
\end{enumerate}

This follows from Prop. 12, (ii).

COROLLARY 3. --- Let w be a common factor of f and g. For w to be a gcd of f and g it is necessary and sufficient that f/w and g/w are relatively prime. This follows from Cor. 2.

COROLLARY 4. --- Let \(f, g, h \in \mathbf{K}[\mathbf{X}]\) . If \(f\) divides \(gh\) and is prime to \(g\) , then \(f\) divides \(h\) .

For \(f\) divides \(gh\) and \(fh\) , hence \(f\) divides every gcd of \(gh\) and \(fh\) , in particular \(h\) (Cor. 2, (i)).

COROLLARY 5. --- Let \(f, g \in \mathbf{K}[\mathbf{X}]\) . For \(f\) and \(g\) to be relatively prime it is necessary and sufficient that the canonical image of \(g\) in \(\mathbf{K}[\mathbf{X}] / (f)\) should be invertible.

For this condition means that there exist \(u, v \in \mathbf{K}[\mathbf{X}]\) such that \(uf + vg = 1\) .

COROLLARY 6. --- Let \(f, g_1, g_2, \ldots, g_n \in \mathbf{K}[\mathbf{X}]\) . If \(f\) is prime to \(g_1, g_2, \ldots, g_n\) , then \(f\) is prime to \(g_1g_2\ldots g_n\) .

* COROLLARY 7. --- For \(f\) and \(g\) to be relatively prime it is necessary and sufficient that they have no common roots in any extension of \(K\) .

For if \(d\) is a gcd of \(f, g\) then the roots common to \(f\) and \(g\) in an extension \(\mathbf{K}'\) of \(\mathbf{K}\) are the roots of \(d\) in \(\mathbf{K}'\) . Now the corollary follows from \(V\) , p.~21, Prop. 4. *

